一位同事把一个函数交给你评审，附言写着：前 1000 个输入都跑过，结果全对。这句话能提高你的信心，却结束不了讨论——第 1001 个输入没跑过，后面还剩无穷多个。测试积累的是证据，积累不出``对所有输入都成立''。

本章要处理的就是这道坎。断言涉及的对象有无穷多个，而人能写下来的论证只有有限行。

出路藏在对象自己的来历里。自然数不是一堆散落的符号，每一个都是从 \(0\) 出发、有限次执行``加一''得到的；一棵语法树不是节点的随意堆积，它由叶子经若干构造规则拼出来；一次循环跑到的每个状态，都是初始状态迭代有限次的结果。既然对象是被有限条规则造出来的，那么论证只要覆盖了每一条规则，就自动覆盖了这些规则能造出的一切。有限的论证于是够用。

\begin{center}
\begin{tikzpicture}[x=1cm,y=1cm,>=stealth,font=\small]
  \node at (0.2,3.0) {\textbf{基础对象}};
  \node at (3.6,3.0) {\textbf{生成规则}};
  \node at (7.6,3.0) {\textbf{相应的证明}};
  \draw[black!30] (-3.0,2.6) -- (9.2,2.6);
  \draw[black!30] (-1.45,-0.5) -- (-1.45,2.6);

  \node[left] at (-1.6,2.0) {自然数};
  \node at (0.2,2.0) {\(0\)};
  \node at (3.6,2.0) {\(n\mapsto n+1\)};
  \node at (7.6,2.0) {数学归纳法};

  \node[left] at (-1.6,1.0) {树与表达式};
  \node at (0.2,1.0) {叶子};
  \node at (3.6,1.0) {由子结构拼装};
  \node at (7.6,1.0) {结构归纳};

  \node[left] at (-1.6,0.0) {循环};
  \node at (0.2,0.0) {初始状态};
  \node at (3.6,0.0) {执行一轮迭代};
  \node at (7.6,0.0) {循环不变量};

  \foreach \y in {0,1,2}{
    \draw[->,black!45] (1.4,\y) -- (2.4,\y);
    \draw[->,black!45] (5.0,\y) -- (6.2,\y);
  }
\end{tikzpicture}

\emph{三类对象各有各的生成方式，而证明方法不过是把生成方式照抄一遍。}
\end{center}

本章四篇讲的是同一个想法落在三种材料上的样子，不是四种彼此独立的技巧。

\subsection{本章篇目}

\begin{itemize}
\tightlist
\item \hyperref[sec:01-Mathematical-Language/04-Induction-and-Recursion/01]{``数学归纳法为何成立''}：基例与归纳步各自防住哪一类失效，少掉任何一个会塌成什么样。
\item \hyperref[sec:01-Mathematical-Language/04-Induction-and-Recursion/02]{``强归纳与良序原理''}：当一步要调用的不止是前一个情形时，归纳假设该开多大；最小反例为什么是同一件事的反面写法。
\item \hyperref[sec:01-Mathematical-Language/04-Induction-and-Recursion/03]{``递归定义与结构归纳''}：怎样用有限条规则钉死一个无限集合，以及封闭条款为什么是结构归纳的合法性来源。
\item \hyperref[sec:01-Mathematical-Language/04-Induction-and-Recursion/04]{``不变量与算法正确性''}：把同一套论证搬到时间方向上，找出进入循环、每一轮、退出时都成立的那句断言。
\end{itemize}

前两篇在自然数上建立原型，第三篇换材料，第四篇换成动态过程。若你是被代码正确性的困惑带进来的，可以直接从第四篇读起，再回头补第一、二篇里关于基例与终止性的讨论。

归纳常被误读成``假设结论成立，再证明结论''。归纳假设针对的只是一个更小的、已被前面论证覆盖过的情形；归纳步要证的是传递机制本身；基例负责让这台机制真的启动。本章反复出现的错误，几乎都能追到这三句话里某一句被悄悄放松了。
